\section{An explicit boundary: FC15 without membership transport}
\label{sec:fc15-e-countermodel}

The $E$-congruence row cannot simply be deleted from the final
20-clause FC15$^*$ target. There is a ZFC-explicit model of all fifteen
native clause templates \emph{with source equality represented by the
predicate $\operatorname{Eq}$}, together with every original completion
row except transport for $E$. In particular, the model satisfies even
all nine function-congruence rows. It does not interpret
$\operatorname{Eq}$ as genuine equality. Thus this is not a model of
the original native FC15$^*$ theory with genuine equality.

\paragraph{Carrier and assignments.}
Take $D=\mathbb N\times\{0,1\}$. Write $|a|$ for the first coordinate
of $a$, put $\widehat n=(n,0)$, and let $\delta=(0,1)$. Interpret
$d=\widehat0$ and
\[
\operatorname{Eq}(a,b)\quad\Longleftrightarrow\quad |a|=|b|.
\]
Let $p_j$ be the $j$th prime, indexed with $p_0=2$, and let
$\nu_p(n)$ denote the exponent of $p$ in the positive integer $n$.
For natural numbers $s,j,z$, define
\[
v(s,j)=\nu_{p_j}(s+1),\qquad
c(s,j,z)=\frac{s+1}{p_j^{v(s,j)}}\,p_j^z-1.
\]
Both functions are total and natural-valued. Unique prime
factorization gives
\[
v(c(s,j,z),j)=z,\qquad
k\ne j\ \Longrightarrow\ v(c(s,j,z),k)=v(s,k).
\]
Set
\[
V(a,b)=\widehat{v(|a|,|b|)},\qquad
C(a,b,z)=\widehat{c(|a|,|b|,|z|)}.
\]
The two displayed identities prove native assignment rows 01 and 02
with the stipulated interpretation of $\operatorname{Eq}$.

\paragraph{Profiles, guard, and witnesses.}
Let $o(n)$ be one when $n$ is odd and zero otherwise. Interpret
\[
\begin{aligned}
\operatorname{Inc}(s,a)&\quad\Longleftrightarrow\quad o(|a|)=1,\\
N(a,b)&=\widehat{1-o(|a|)o(|b|)},&
A(j,k)&=\widehat0,\\
Q(j,a)&=\widehat{|a|},&
W(s,j,a)&=\delta,\\
H(b,t)&=\widehat0,&
Y(b,t,u)&=Z(b,t,u)=(|u|+1,1).
\end{aligned}
\]
The membership relation is the diagonal on decoration one alone:
\[
E((n,i),(m,j))\quad\Longleftrightarrow\quad i=j=1\ \text{and}\ n=m.
\]
Finally set
\[
G(b,t)\quad\Longleftrightarrow\quad |t|\text{ is even}.
\]
The three NAND rows 03--05 hold because incidence with $N(a,b)$ is
exactly the NAND of incidence with $a$ and $b$. Every $V$ value has
second coordinate zero. Hence $E(V(s,j),V(s,k))$ is always false,
as is $\operatorname{Inc}(s,A(j,k))$; this proves rows 06 and 07.
Incidence ignores its first argument and $Q$ preserves the parity of
its second argument, so rows 08 and 09 hold. The noncanonical constant
value of $W$ causes no failure of row 09.

\paragraph{All six closure rows.}
For the exact relation abbreviation in the native templates,
\[
R(b,t,x,y)=\operatorname{Inc}(C(C(b,d,x),q,y),t),
\]
we have $R(b,t,x,y)$ if and only if $|t|$ is odd, independently of
$b,x,y$. If $|t|$ is odd, $G$ is false and $R$ is full. Row 10 then
holds through $R$. Row 11 holds because $Y$ and $Z$ are identical
elements on the decoration-one diagonal. Row 12 holds because the
first coordinate of $Z(b,t,u)$ is $|u|+1$, so $E(Z(b,t,u),u)$ is
false. Rows 13--15 are true through $\neg G$.

If $|t|$ is even, $G$ is true and $R$ is empty. Rows 10--12 hold
through $G$. Rows 13 and 14 hold through $\neg R$, regardless of the
derived-member term $D^*$ in row 14. Row 15 holds because $H(b,t)$
has decoration zero and therefore has no incoming $E$-edge.
This proves all fifteen native templates for arbitrary elements of
the infinite carrier, including every colliding variable assignment.

\paragraph{Completion and a ground failure of $E$ transport.}
The predicate $\operatorname{Eq}$ is an equivalence relation.
Every function above has output first coordinate determined solely by
the first coordinates of its inputs. Coordinatewise
$\operatorname{Eq}$-equal arguments therefore give
$\operatorname{Eq}$-equal outputs, proving all nine
function-congruence rows. Incidence and the guard depend only on the
first coordinate of their second argument, so both predicate transport
rows hold. The selector $q=N(d,d)$ is $\widehat1$ and is not
$\operatorname{Eq}$-equivalent to $d$.

But the ground term $w=W(d,d,d)$ denotes $\delta$. Thus
\[
\operatorname{Eq}(w,d),\qquad E(w,w),\qquad \neg E(d,d)
\]
all hold. These falsify the $E$-transport instance moving both arguments
from $w$ to $d$. The failure is witnessed by ground terms, not merely
by unnamed elements of a larger structure. The whole construction is
elementary arithmetic formalizable in ZFC. Consequently the
19-clause FC15$^*$ predicate theory obtained by simply removing $E$
transport from the final 20-clause target has a ZFC-provable model.
This rules out that particular further deletion as a route to the
required consistency contract; it does not rule out a different
presentation or reconstruction with other axioms.

\paragraph{The FC16$^*$ limitation is exact.}
This model does not satisfy the additional coded-extensionality row.
Choose $a=d=\widehat0$, $b=q=\widehat1$, and
$j=C(d,q,q)=\widehat2$. All three equality guards are false. Both
atomic class values $A(j,d),A(j,q)$ are $\widehat0$, their NAND
exclusive-or class is $\widehat0$, and $Q$ preserves that value.
The final incidence literal is therefore false as well. The model
settles the simple deletion for FC15$^*$ only.

\paragraph{Executable diagnostic.}
The companion \texttt{fc15\_without\_e\_countermodel.py} independently
evaluates the immutable native clause trees and all retained completion
trees on every assignment from
$\{0,1,2\}\times\{0,1\}$. Intermediate function values are evaluated
as unbounded integers; the sample is not a finite carrier. It also
checks the exact ground $E$-transport failure and the added FC16$^*$
row failure. These finite evaluations are diagnostics of the displayed
infinite-model proof, not a substitute for that proof.
